%% -*- coding: utf-8 -*-
\documentclass[12pt,pagesize,paper=192mm:108mm]{scrbook} 
%1920x1080 1280x720
\areaset[current]{192mm}{108mm}
\usepackage{calc}
\usepackage[T2A]{fontenc}
\usepackage[utf8]{inputenc}
\usepackage[english,russian]{babel}
\usepackage{microtype}
\usepackage{misccorr}
\usepackage{cmap}
%\usepackage[unicode=true]{hyperref}
\usepackage{graphicx}
\usepackage{amssymb}
\usepackage{amsmath}
%\usepackage{srcltx}
\usepackage{textcomp}
\usepackage{xspace}
%научные символы и смайлики \smiley \frownie
\usepackage{wasysym}
\usepackage{ccicons}
\begin{document}
\begin{titlepage}
  \vspace*{-0.5em}
  \begin{center}    
    \hspace*{3em}
    \begin{minipage}[t]{3em}
      \includegraphics[width=\textwidth]{../BINP-logo}
    \end{minipage}\hfill
    \begin{minipage}{0.23\linewidth}
    \includegraphics[width=\textwidth]{../NSU-logo}
    \end{minipage}
    \hfill
    \hspace*{6em}

    Кафедра теоретической физики физического факультета НГУ
    \medskip

    \Large
    Профессор Фадин В.\,С.
    \bigskip

    \huge
    \textbf{Квантовая электродинамика}
    \bigskip

    \Large
    Лекция № 4
    \vfill

    \normalsize
    % \begin{minipage}{0.65\linewidth}
    % \end{minipage}
    \vfill

    \normalsize \ccbysa\hspace{0.5em}  Новосибирск 2013
  \end{center}
\end{titlepage}
\newpage

\vspace*{-1em}
\begin{center}
\vfill
  \begin{minipage}{0.75\linewidth}
    Простейшие процессы КЭД. Рассеяние электрона во внешнем поле:
    диаграмма Фейнмана, матричный элемент, дифференциальное по углу
    вылета конечного электрона сечение рассеяния, нерелятивистский
    предел, ультрарелятивистский предел (подавление рассеяния назад и
    сохранение спиральности), расходимость полного сечения. Рассеяние
    электрона на скалярной частице $e^-\pi^-\to e^-\pi^-$: правила
    Фейнмана для взаимодействия скалярных частиц с электромагнитным
    полем из принципа калибровочной инвариантности и принципа
    соответствия в нерелятивистском пределе, матричный элемент, учет
    внутренней структуры скалярной частицы (пиона) или распределения
    заряда рассеивающего центра (формфактор), рассеяние назад,
    дифференциальное сечение в системе центра инерции,
    мандельстамовские переменные. Кросс"=инвариантность: квадрат
    матричного элемента для процесса $e^+e^-\to\pi^+\pi^-$,
    дифференциальное сечение в системе центра инерции, проявление
    закона сохранения спиральности в канале аннигиляции, зависимость
    сечения от скорости конечных пионов.
  \end{minipage}
  \vfill

  % \normalsize \ccbysa\hspace{0.5em} Новосибирск 2013
\end{center}
\end{document}
